%! Regression Lines and Residuals — Unit 2, Lesson 2.3 — Guided Notes & Practice (Answer Key)
% THE in-class document, in the MAIN IDEAS / NOTES shape (LESSON_SHAPE.md §1), at the COMPACT
% budget: the vocabulary box, then ONE two-column guidednotes table of four rows, 14 problems,
% two pages. Problems are numbered 1..N through the whole component.
%   I DO   : the teacher fills the definition lines and works the FIRST problem of each grid
%            (1, 4, 7, 11).
%   WE DO  : the class works the rest of each grid, students holding the pen; the last row,
%            Guided Practice, is worked entirely together in a fresh context.
%   YOU DO : nothing on this page — ../homework/ IS the individual practice.
% THE TRAP is problem 3 (which number is y and which is y-hat). THE CRUX is problem 9, in the
% last instruction row: a residual computed as y-hat minus y, and a positive residual read as
% "the line was too high" (EK DAT-1.E.1).
%
% THE DATA (instruction rows): 17 pizza deliveries, 1 to 9 miles; line time-hat = 8 + 3(distance).
%   (x, residual): (1,1) (1.5,-2) (2,2) (2.5,-1) (3,0) (3.5,2) (4,-3) (4.5,1) (5,3) (5.5,-2)
%                  (6,-2) (6.5,1) (7,-1) (7.5,2) (8,0) (8.5,-2) (9,1)   — sum 0, no pattern.
% GUIDED PRACTICE: 20 adults, 4 to 9 hours of sleep; RT-hat = 430 - 25(sleep); the 7-hour
%   adult at 270 ms has residual +15.
%
% PAGE LOCKSTEP with notes_key/: every \blank{} becomes an \ans{}, every \termblank a \termans,
% and every \pcell gets its answer in the fourth argument. Nothing else differs.
% Slope/intercept INTERPRETATION and r^2 are Lesson 2.4 — the equation is only USED here.
% NO SKETCH QUESTIONS: the scatterplot and both residual plots are pre-drawn in TikZ.
\documentclass[12pt]{article}
\usepackage{apstats-article}
\usepackage{apstats-key}   % pulls in -boxes; do NOT also load -boxes

\begin{document}

\pageheader{Unit 2, Lesson 2.3}{Guided Notes \& Practice --- Answer Key}
% NAMESTRIP: no name row here — it lives on the cover only.

\vspace{2pt}

\boxguard[16]
\begin{vocabbox}
\small
Fill in each term \textbf{as we name it} in the notes below.
\par
\termans{Regression line}{a line $\hat{y}=a+bx$ that uses the explanatory variable $x$ to predict the response variable $y$.}
\par
\termans{Predicted value ($\hat{y}$)}{the $y$-value the line gives for a particular $x$; the actual $y$ usually differs from it.}
\par
\termans{Extrapolation}{predicting for an $x$ outside the range of $x$-values used to build the line --- unreliable.}
\par
\termans{Residual}{actual minus predicted, $y-\hat{y}$; positive means the point lies above the line.}
\par
\end{vocabbox}

\small
\begin{guidednotes}
% ── ROW 1: the equation and prediction ───────────────────────────────────────
\mainidea[The equation \&]{Prediction} &
A \textbf{regression line} $\hat{y}=a+bx$ models how the \ans{response} variable $y$ changes
as the \ans{explanatory} variable $x$ changes. $\hat{y}$ (``$y$-hat'') is the \ans{predicted}
value; $y$ is the \ans{actual} value. To predict, \ans{substitute} $x$ into the equation.

\begin{minipage}[c]{0.60\linewidth}\raggedright
A pizza shop timed 17 deliveries, 1 to 9 miles out (time in minutes, distance in miles):\par
\vspace{3pt}
\centering $\widehat{\text{time}} = 8 + 3(\text{distance})$\par
\end{minipage}\hfill
\begin{minipage}[c]{0.38\linewidth}
\centering
\begin{tikzpicture}[xscale=0.34, yscale=0.048]
  \foreach \y in {10,20,30,40} {\draw[linegray, very thin] (0,\y) -- (10,\y);
    \node[left, font=\tiny] at (0,\y) {\y};}
  \draw[->, thick] (0,0) -- (10.4,0);
  \draw[->, thick] (0,0) -- (0,43);
  \foreach \x in {0,2,4,6,8,10} {\draw (\x,-0.8) -- (\x,0.8);
    \node[below, font=\tiny] at (\x,-0.5) {\x};}
  \draw[goldacc, very thick] (1,11) -- (9,35);
  \foreach \x/\y in {1/12, 1.5/11.5, 2/16, 2.5/14.5, 3/17, 3.5/20.5, 4/17, 4.5/22.5, 5/26,
                     5.5/22.5, 6/24, 6.5/28.5, 7/28, 7.5/32.5, 8/32, 8.5/31.5, 9/36}
    {\fill[navy] (\x,\y) circle[x radius=0.15, y radius=1.06];}
  \node[font=\tiny] at (5,-11.5) {Distance (miles)};
  \node[rotate=90, font=\tiny] at (-2.3,21) {Time (min)};
\end{tikzpicture}
\end{minipage}
\notesprompt{Use the line.}
\begin{probgrid}{|Y|Y|Y|}
\pcell{1}{Predict the time for a 5-mile delivery.}{1.0cm}{$8+3(5)=23$ minutes} &
\pcell{2}{Predict the time for a 2.5-mile delivery.}{1.0cm}{$8+3(2.5)=15.5$ minutes} &
\pcell{3}{A 6-mile delivery took 24 minutes. Give $y$ and $\hat{y}$.}{1.0cm}{$y=24$ (actual), $\hat{y}=26$ (predicted)} \\ \hline
\end{probgrid}
\\ \hline
% ── ROW 2: extrapolation ─────────────────────────────────────────────────────
\mainidea[Predicting beyond]{The data} &
\textbf{Extrapolation} is predicting for an $x$-value \ans{outside} the interval of $x$-values
used to build the line --- the farther out, the \ans{less} reliable. Say so every time.
\notesprompt{The shop's line was built from deliveries of 1 to 9 miles.}
\begin{probgrid}{|Y|Y|Y|}
\pcell{4}{Predict the time for a 30-mile delivery. Trust it?}{1.2cm}{$98$ min --- no: far outside 1--9 mi} &
\pcell{5}{Which are extrapolation: 3 miles, 8.5 miles, 12 miles?}{1.2cm}{Only 12 miles (outside 1--9)} &
\pcell{6}{Could the line be wrong for 15 miles? Why can't we check?}{1.2cm}{Yes --- no data past 9 mi to check} \\ \hline
\end{probgrid}
\\ \hline
% ── ROW 3: residuals — THE CRUX ──────────────────────────────────────────────
\mainidea[Actual minus]{Predicted} &
A \textbf{residual} is actual minus predicted: \quad
residual $=$ \ans{$y$} $-$ \ans{$\hat{y}$}

\vspace{2pt}
\begin{minipage}[c]{0.60\linewidth}\raggedright
\textbf{Positive:} the actual value is \ans{larger} than predicted --- the point sits above
the line, so the line predicted too \ans{low}.

\vspace{3pt}
A \textbf{residual plot} graphs each residual against $x$. Points scattered \ans{randomly}
around 0, with no pattern, say a linear model is \ans{appropriate}.
\end{minipage}\hfill
\begin{minipage}[c]{0.38\linewidth}
\centering
\begin{tikzpicture}[xscale=0.34, yscale=0.30]
  \draw[->, thick] (0,-4.6) -- (0,4.8);
  \draw[thick] (0,-4.6) -- (10.2,-4.6);
  \draw[dashed, slate] (0,0) -- (10.2,0);
  \foreach \y in {-4,-2,0,2,4} {\draw (-0.12,\y) -- (0.12,\y);
    \node[left, font=\tiny] at (0,\y) {\y};}
  \foreach \x in {0,2,4,6,8,10} {\draw (\x,-4.75) -- (\x,-4.45);
    \node[below, font=\tiny] at (\x,-4.6) {\x};}
  \foreach \x/\e in {1/1, 1.5/-2, 2/2, 2.5/-1, 3/0, 3.5/2, 4/-3, 4.5/1, 5/3, 5.5/-2,
                     6/-2, 6.5/1, 7/-1, 7.5/2, 8/0, 8.5/-2, 9/1}
    {\fill[navy] (\x,\e) circle[x radius=0.15, y radius=0.17];}
  \node[font=\tiny] at (5,-6.2) {Distance (miles)};
  \node[rotate=90, font=\tiny] at (-2.1,0) {Residual (min)};
\end{tikzpicture}
\end{minipage}
\\
& \notesprompt{Residuals for the pizza deliveries.}
\begin{probgrid}{|Y|Y|}
\pcell{7}{The 5-mile delivery took 26 minutes. Find its residual.}{1.4cm}{$26-23=+3$ minutes} &
\pcell{8}{A 4-mile delivery took 17 minutes. Find its residual and say what it means.}{1.4cm}{$17-20=-3$: faster than predicted; the line overpredicted} \\ \hline
\pcell{9}{Jordan says the 5-mile residual is $23-26=-3$, ``so the line predicted too high.''
Fix both mistakes.}{1.5cm}{$y-\hat{y}=26-23=+3$. It took longer than predicted: the line predicted too \emph{low}.} &
\pcell{10}{Does the residual plot show a pattern? Is a linear model appropriate?}{1.5cm}{No pattern --- random scatter around 0. Yes, linear is appropriate.} \\ \hline
\end{probgrid}
\\ \hline
% ── ROW 4: guided practice — WE DO, together, fresh context ──────────────────
\mainidea[Guided practice]{Sleep \& Reaction} &
\begin{minipage}[c]{0.60\linewidth}\raggedright
A researcher measured hours of sleep and reaction time, in milliseconds, for 20 adults who
slept 4 to 9 hours:\par
\centering $\widehat{\text{RT}} = 430 - 25(\text{sleep})$\par
\raggedright
Its residual plot is at right.
\end{minipage}\hfill
\begin{minipage}[c]{0.38\linewidth}
\centering
\begin{tikzpicture}[xscale=0.55, yscale=0.050]
  \draw[->, thick] (3.5,-27) -- (3.5,28);
  \draw[thick] (3.5,-27) -- (9.5,-27);
  \draw[dashed, slate] (3.5,0) -- (9.5,0);
  \foreach \y in {-20,0,20} {\draw (3.42,\y) -- (3.58,\y);
    \node[left, font=\tiny] at (3.5,\y) {\y};}
  \foreach \x in {4,5,6,7,8,9} {\draw (\x,-28) -- (\x,-26);
    \node[below, font=\tiny] at (\x,-27) {\x};}
  \foreach \x/\e in {4/10, 4.5/-15, 4.5/5, 5/-20, 5/15, 5.5/0, 6/20, 6/-10, 6.5/-5, 6.5/12,
                     7/15, 7/-18, 7.5/5, 7.5/-8, 8/-12, 8/18, 8.5/3, 8.5/-15, 9/10, 9/-10}
    {\fill[navy] (\x,\e) circle[x radius=0.09, y radius=1.0];}
  \node[font=\tiny] at (6.5,-37) {Sleep (hours)};
  \node[rotate=90, font=\tiny] at (2.35,0) {Residual (ms)};
\end{tikzpicture}
\end{minipage}
\\
& \notesprompt{We work these together.}
\begin{probgrid}{|Y|Y|}
\pcell{11}{Predict the reaction time after 7 hours of sleep.}{1.4cm}{$430-25(7)=255$ ms} &
\pcell{12}{An adult who slept 7 hours reacted in 270 ms. Find the residual and say what it
means.}{1.4cm}{$270-255=+15$ ms: slower than predicted; the line underpredicted} \\ \hline
\pcell{13}{Predict the reaction time after 12 hours of sleep. Trust it?}{1.5cm}{$130$ ms. No --- 12 h is outside 4--9 h (extrapolation)} &
\pcell{14}{Is a linear model appropriate here? Use the residual plot.}{1.5cm}{Yes: the residuals scatter randomly around 0, no pattern} \\ \hline
\end{probgrid}
\\ \hline
\end{guidednotes}

% ── THE NOTES END AT GUIDED PRACTICE ─────────────────────────────────────────
% There is deliberately no "you do" here: ../homework/ is the individual practice,
% started in class in the last 8 minutes while the teacher circulates.

\end{document}
